% SAI / IEEE-style two-column conference format (IEEEtran.cls is not installed
% locally, so we emulate the SAI_PAPER_FORMAT template in the article class:
% two columns, US-letter, Roman-numeral section heads, "Abstract--"/"Keywords--"
% run-ins, and numeric references in citation order).
%
% Multi-model, reviewer-revised edition. Reuses the preamble, macros, and thesis
% of paper.tex; rebuilds the empirical section from the three-model, four-run-set
% run set (llama3.1:8b, qwen3:4b x2 independent run sets, qwen2.5:32b,
% 2026-08-11); adds the delegation
% scenario, the discovery / tool-prompt coupling levers, a reproducible technical
% specification of GCP, and a two-baseline framing of the provenance metric.
\documentclass[10pt,twocolumn,letterpaper]{article}

% --- packages (all confirmed present in the local TeX install) ---
\usepackage[utf8]{inputenc}
\usepackage[T1]{fontenc}
\usepackage{lmodern} % scalable fonts, required for microtype font expansion
\usepackage[letterpaper,top=0.75in,bottom=1in,left=0.625in,right=0.625in,columnsep=0.25in]{geometry}
\usepackage{amsmath,amssymb}
\usepackage{graphicx}
\usepackage{booktabs}
\usepackage{multirow}
\usepackage{caption}
\usepackage{microtype}
\usepackage{xcolor}
\usepackage{verbatim}
\usepackage[numbers,sort&compress]{natbib}
\usepackage{xurl} % allow long reference URLs to break across lines
\usepackage[hidelinks]{hyperref}

% IEEE/SAI section numbering: Roman sections (I, II, ...), lettered subsections.
\renewcommand{\thesection}{\Roman{section}}
\renewcommand{\thesubsection}{\Alph{subsection}}
\renewcommand{\thesubsubsection}{\arabic{subsubsection}}
\captionsetup{font=small,labelsep=period}

\newcommand{\gcp}{\textsc{gcp}}
\newcommand{\ata}{\textsc{a2a}}
\newcommand{\bigoh}{\mathcal{O}}

\title{\textbf{Read-First, Role-Gated Context Federation for Multi-Agent LLM
Systems:\\ A Fairness-Controlled, Multi-Model Evaluation Against\\
Point-to-Point Message-Passing}}

% IEEE-style author block: fixed-width author columns laid out side by side,
% name / affiliation (italic) / city, country / email. Three authors fit the
% full-width title band produced by \maketitle in twocolumn mode.
\author{%
\begin{minipage}[t]{0.32\textwidth}\centering
Adrian Issac Mamani Quispe\\
\textit{Univ.\ Nacional de San Agust\'in}\\
\textit{de Arequipa}\\
Arequipa, Peru\\
admamaniq@unsa.edu.pe
\end{minipage}%
\hfill
\begin{minipage}[t]{0.32\textwidth}\centering
Ryan Fabian Valdivia Segovia\\
\textit{Univ.\ Nacional de San Agust\'in}\\
\textit{de Arequipa}\\
Arequipa, Peru\\
rvaldiviase@unsa.edu.pe
\end{minipage}%
\hfill
\begin{minipage}[t]{0.32\textwidth}\centering
Karim Guevara Puente de la Vega\\
\textit{Univ.\ Nacional de San Agust\'in}\\
\textit{de Arequipa}\\
Arequipa, Peru\\
kguevarap@unsa.edu.pe
\end{minipage}%
}

\date{August 2026}

\begin{document}
\maketitle

\begin{abstract}
\textbf{Problem.} Multi-agent systems based on large language models (LLMs) commonly rely on
point-to-point message-passing for agent coordination, which can increase integration complexity
and requires additional mechanisms for fine-grained access control and traceability.
\textbf{Method.} This study evaluates the Graph Context Protocol (\gcp{}), a read-first
context-federation approach that integrates role- and capability-based access control with per-read
provenance. \gcp{} is compared with a base Agent2Agent (\ata{}) message-passing implementation
under a controlled experimental design in which both approaches use the same agent logic, models,
tasks, and evaluation criteria, differing only in the interoperability layer. The evaluation
comprises four scenarios (marketplace, software organization, cross-owner supply chain, delegation)
and three open-weight LLMs (\texttt{llama3.1:8b}, \texttt{qwen3:4b}, \texttt{qwen2.5:32b}) across
four independent run sets---\texttt{qwen3:4b} is evaluated in two independent run sets, not a fourth
distinct model---3--5 seeds per arm, and distinguishes deterministic architectural properties from
model-dependent behavioral outcomes. \textbf{Results.} Structural analysis shows that \gcp{} reduces
integration-edge growth from $\bigoh(n^2)$ to $\bigoh(n)$ ($50$ vs.\ $2450$ at $n{=}50$) and
maintains constant per-request discovery and tool-description overhead as the number of peers
increases, while both grow linearly for base \ata{}. In the delegation scenario, \gcp{} prevented
all attempted unauthorized actions (containment $1.00$), whereas base \ata{} executed unauthorized
actions whenever the model attempted them. \gcp{} also provided complete provenance by construction,
while equivalent per-read traceability in \ata{} requires additional instrumentation. Behavioral
experiments showed comparable task success and communication cost between approaches, while current
canary-based leakage experiments showed no behavioral privacy advantage for either architecture.
These results indicate that \gcp{} provides model-independent structural and governance benefits
while maintaining comparable behavioral performance. \textbf{Limitations.} Behavioral results use
open-weight models and 3--5 seeds rather than the 20--30 recommended for significance testing on the
token/latency/success deltas; this affects only those preliminary behavioral columns, not the
deterministic/architectural results (coupling, discovery, tool-prompt, provenance, containment).
\textbf{Contribution.} The contribution lies in integrating federated read-first access, role- and
capability-based control, and native per-read provenance within a unified multi-agent coordination
substrate, supported by a fairness-controlled, multi-model evaluation.
\end{abstract}

\smallskip
\noindent\textbf{\textit{Keywords---}}\textit{context federation; multi-agent
systems; large language models; access control; least privilege; provenance;
agent containment; software coupling; agent communication protocols}
\smallskip

\section{Introduction}

The dominant coordination pattern for LLM-based multi-agent systems is direct message-passing.
Frameworks such as AutoGen~\citep{wu2023autogen}, ChatDev~\citep{qian2024chatdev}, and
MetaGPT~\citep{hong2024metagpt} compose applications from agents that send and receive messages,
and the Agent2Agent (A2A) protocol~\citep{a2a2025} standardizes exactly this: a client agent
advertises and delegates tasks to named remote agents over JSON-RPC. This paradigm descends from
the classical agent communication languages KQML~\citep{finin1994} and FIPA-ACL~\citep{fipa2002},
which already assumed pairwise connections and shared, pre-negotiated protocols between every
interacting pair.

Message-passing has a structural cost the software-engineering literature has understood for fifty
years. Connecting $n$ components pairwise requires on the order of $n^2$ links; a shared
intermediary that each component connects to once reduces this to $n$. This is the ``integration
spaghetti'' the Enterprise Integration Patterns~\citep{hohpe2003} catalog addresses with brokers
and hubs, and the phenomenon coupling theory~\citep{parnas1972,stevens1974,briand1999} formalizes:
modifiability is governed by the number and strength of inter-component dependencies.

We study an alternative substrate, the Graph Context Protocol (\gcp{}). \gcp{} is
\emph{read-first}: an agent obtains what it needs by reading context from independently owned
knowledge nodes rather than receiving pushed messages. It is \emph{role-gated}: every read passes
through a policy gate that authorizes the request by role and required capabilities, defaulting to
deny~\citep{saltzer1975,sandhu1996}. And it is \emph{audited}: every read or delegation
decision---grant or deny---produces a structured provenance record.

\paragraph{Thesis (measured form).} At equal task success, read-first role-gated context federation
offers \emph{structural} advantages over point-to-point message-passing in coupling and
traceability, and \emph{preliminary} evidence of lower or equal cost under certain scenarios. We
state the structural claims as strong (they are deterministic or architectural) and the behavioral
claims as preliminary (open-weight models, few seeds).

\paragraph{Scope of the comparison.} We are explicit about what the A2A arm represents. We do not
claim message-passing \emph{cannot} be audited: logs, middleware, proxies, or gateways can add an
access trail. Our comparison is against \emph{base} A2A---the protocol as shipped, without such
instrumentation---and its purpose is to show that \gcp{} makes provenance and role gating
\emph{native primitives} rather than optional add-ons. Where useful we discuss a second reference
point, A2A augmented with minimal read logging (Section~\ref{sec:metrics}).

\paragraph{Contributions.}
\begin{enumerate}
  \item \textbf{Architecture and a reproducible specification.} We describe \gcp{} (read-first,
  role-gated, audited, with a stricter-capability delegation gate) precisely enough to reimplement:
  query format, node and policy structure, role credential, provenance record, and the
  authorization algorithm (Section~\ref{sec:arch}).
  \item \textbf{A fairness-controlled benchmark method.} One scenario definition and one injected
  model build both arms; only the interop layer changes (Section~\ref{sec:method}).
  \item \textbf{A multi-model empirical head-to-head} over four scenarios and three open-weight
  models across four independent run sets, separating deterministic from model-dependent results
  (Section~\ref{sec:results}).
  \item \textbf{Metric definitions} for structural coupling, discovery/tool-prompt overhead, a
  four-way leakage/agency taxonomy, provenance completeness, and delegation containment
  (Section~\ref{sec:metrics}).
\end{enumerate}

\section{Background and Related Work}
\label{sec:related}

\subsection{Message-passing as the multi-agent baseline}
Agent communication languages established the message-passing paradigm: KQML~\citep{finin1994}
wrapped payloads in typed performatives over pairwise connections, and FIPA-ACL~\citep{fipa2002}
standardized a message in which sender and receiver must share content language and ontology before
any exchange succeeds. Modern LLM frameworks reinvented this in natural language.
AutoGen~\citep{wu2023autogen} programs applications as conversations; ChatDev~\citep{qian2024chatdev}
arranges role-pairs in a fixed ``chat chain''; surveys~\citep{guo2024survey} taxonomize communication
into peer-to-peer, centralized, layered, and shared-pool variants. The Model Context
Protocol~\citep{mcp2024} standardizes agent-to-tool context access (an early read-first form), while
A2A~\citep{a2a2025} standardizes agent-to-agent task delegation---the baseline we benchmark against.
\citet{zhang2024agentprune} report 28--73\% excess tokens in inter-agent communication graphs and
recover most by pruning topology; \citet{mao2026delm} replace both centralized orchestration and
peer-to-peer routing with a shared verified context at roughly half the cost. We add role-gated
authorization, audited per-read provenance, an agency-containment gate, and a multi-model
fairness-controlled comparison.

\subsection{Coupling and the $n^2$ problem}
That fewer inter-component links yield more maintainable systems is among the oldest results in
software engineering. \citet{parnas1972} showed information-hiding decomposition produces far lower
coupling than flowchart-style partitioning; \citet{stevens1974} made coupling and cohesion
measurable; \citet{briand1999} gave a graph-theoretic framework in which coupling is monotone in
dependency-edge count, so an $\bigoh(n)$-edge topology is provably less coupled than an
$\bigoh(n^2)$ one. \citet{hohpe2003} translated this to integration: a point-to-point mesh of $n$
systems needs $n(n{-}1)/2$ channels, which a hub reduces to $n$. Microservice
research~\citep{panichella2021} carries the same argument into cloud-native systems, and
\citet{shen2025topology} show dense agent topologies amplify error propagation while sparser ones
suppress it. \ata{} creates agent$\leftrightarrow$agent edges; \gcp{} creates agent$\rightarrow$store
edges.

\subsection{Least privilege, access-controlled context, and gated retrieval}
\gcp{}'s read gate operationalizes least privilege~\citep{saltzer1975}. Role-based access
control~\citep{sandhu1996} assigns permissions to roles; attribute-based access
control~\citep{nist800162} generalizes to arbitrary attributes; capability-based
protection~\citep{dennis1966} ties authority to explicitly passed tokens---the model \gcp{}'s role
credentials follow. Data-minimization law~\citep{gdpr2016} makes ``read only what your role
requires'' a legal constraint. Read-first tool access (MCP) and access-controlled retrieval supply
the read side, and RBAC/ABAC the gate side, but neither couples the two with per-read provenance;
the point-to-point baseline has no analogous gate, so leakage and delegation control rest on
sender-side discretion.

\subsection{Leakage, provenance, and governance in agentic systems}
Confidentiality in agentic systems is a system property, not only a model
property~\citep{evertz2024whispers}; inter-agent message channels are the dominant leakage vector,
largely invisible to output-only auditing~\citep{elyagoubi2026agentleak}; and restricting what an
agent can \emph{see} is more robust than filtering what it can \emph{say}~\citep{bagdasarian2024airgap}.
The OWASP Top 10 for LLM applications~\citep{owasp2025} names Sensitive Information Disclosure and
Excessive Agency as primary risks---the latter is exactly what our delegation scenario probes.
Deterministic privilege control at the tool boundary~\citep{shi2025progent}, authenticated
delegation of authority~\citep{south2025delegation}, and audit trails~\citep{ojewale2025audit} argue
for tamper-evident records of system actions.

\paragraph{Delta / defensible novelty.} Prior work supplies the pieces---read-first access (MCP),
role/attribute gating (RBAC/ABAC, access-controlled RAG), audited delegation, and audit logs---but
not their conjunction, and not a fairness-controlled, multi-model head-to-head against a real
message-passing baseline. Our claim to novelty is precisely the \emph{integration}: federated read
$+$ role control $+$ per-read provenance, evaluated against \ata{}. Read-first alone, RBAC alone, or
auditing alone would not be a contribution.

\section{The Graph Context Protocol: Specification}
\label{sec:arch}

\gcp{} models an ecosystem as independently owned \emph{knowledge nodes} and \emph{agent nodes}.
Coordination is read-first: an agent issues a \emph{context query} to a peer node, which returns
context the caller is authorized to read. We specify the substrate precisely enough to reimplement.

\subsection{Data structures}
A \textbf{knowledge node} carries an id, a kind, an adapter that produces content, and an access
policy. A \textbf{policy} declares gating fields; a \textbf{role credential} names the caller's role
and flat capabilities; a \textbf{provenance record} logs each decision.

{\small
\begin{verbatim}
KnowledgeNode {
  id, kind,               // e.g. "rag", "log"
  adapter,                // yields raw content
  policy: AccessPolicy }
AccessPolicy {
  readableByRoles: [roleId],
  requiredCapabilities: [capId],
  fallbackAllowed: bool,
  denialMode: "empty" | "error" }
Principal (role credential) {
  id, role: { id, capabilities:[capId],
              parent? },
  capabilities: [capId] }
Provenance {
  principalId, targetNodeId, timestamp,
  decision: "grant" | "deny",
  matchedPolicy }
\end{verbatim}}

\subsection{Query and authorization flow}
A context query names a target node and a requested capability (read, or the stricter delegate).
The server evaluates the admission predicate \emph{before} any adapter runs, so denied content
never enters the agent's context window (complete mediation~\citep{saltzer1975}). The effective
capability set of a principal is the union of its flat capabilities and those inherited through the
role chain (own-first precedence). The predicate is
\begin{align*}
\mathit{authorized}(P,K)\equiv\;
&\underbrace{(P.\textit{roles}{=}[\,]\vee K.\textit{role}{\in}P.\textit{roles})}_{\textit{roleOk}}\\
\wedge\;
&\underbrace{(P.\textit{caps}{=}[\,]\vee P.\textit{caps}{\subseteq}\mathit{eff}(K))}_{\textit{capsOk}}
\end{align*}
An absent or schema-invalid policy is treated as a denial (default-deny). The end-to-end path is:

\begin{center}\small
Request $\rightarrow$ Policy check $\rightarrow$ Decision
$\rightarrow$ Context returned / denied $\rightarrow$ Provenance record
\end{center}

\subsection{Read/authorize algorithm (pseudocode)}
{\footnotesize
\begin{verbatim}
handleQuery(query, principal):
  node = resolve(query.targetNodeId)
  if node == null: return deny("no-node")
  P = node.policy
  if P invalid or absent:        # default-deny
      record(principal,node,"deny","none")
      return deny("no-policy")
  eff = flat(principal) UNION
        rolechain(principal.role)
  roleOk = P.roles==[] or
           principal.role.id in P.roles
  capsOk = P.caps==[] or P.caps subset eff
  needCap = query.capability     # read|delegate
  capOk2 = needCap in eff
  if roleOk and capsOk and capOk2:
      record(principal,node,"grant",P.id)
      return node.adapter.read()  # content
  else:
      record(principal,node,"deny",P.id)
      return denialResult(P.denialMode)  # denied
\end{verbatim}}
Delegation reuses this path with \texttt{needCap = delegate}, a strictly stronger capability than
\texttt{read}: an agent authorized to read is not thereby authorized to command.

\subsection{Denial handling}
On denial the handler returns only a denial reason (or an empty result, per \texttt{denialMode});
the adapter is never invoked, so node content cannot appear in the response. Every denial is itself
recorded, yielding a who-tried-what ledger, not only a who-read-what one.

\subsection{Relation to A2A and MCP}
Table~\ref{tab:props} contrasts the three. A2A is agent-to-agent and push-based, with coarse
card-declared allow/deny and no native per-read trail; MCP is read-first but agent-to-tool and
single-owner, without role gating across owners; \gcp{} combines cross-owner read-first access,
role gating, deny-by-default, and native per-read provenance.

\begin{table}[t]
\centering
\caption{Property comparison. A ``no'' does not mean the protocol makes the property impossible;
every cell below is achievable in principle by extending the base protocol. ``not nat.'' $=$ not a
native primitive of the base protocol (would require added logic on top of it); ``ext.'' $=$
achievable only with additional external instrumentation; ``n/a'' $=$ the concept does not apply to
this protocol's design.}
\label{tab:props}
\setlength{\tabcolsep}{4pt}
\begin{tabular}{lccc}
\toprule
property & \ata{} & MCP & \gcp{} \\
\midrule
agent$\leftrightarrow$agent            & yes & not nat. & yes \\
read-first                             & not nat. & yes & yes \\
cross-owner boundary                   & yes & not nat. & yes \\
role/capability gating                 & coarse & not nat. & yes \\
per-read provenance (native)           & ext. & ext. & yes \\
deny-by-default                        & not nat. & not nat. & yes \\
stricter-capability delegation         & not nat. & n/a & yes \\
\bottomrule
\end{tabular}
\end{table}

\section{Metric Definitions}
\label{sec:metrics}

We report metrics in four families and separate deterministic from model-dependent throughout.

\paragraph{Structural coupling (deterministic).} For $n$ peers, \emph{pairwise connections} are the
integration edges the topology requires: $C_{\ata{}}(n)=n(n{-}1)$ vs.\ $C_{\gcp{}}(n)=n$.
\emph{Integration effort} is the marginal cost $C(n{+}1)-C(n)$: constant $1$ for \gcp{}, $2n$ for
\ata{}. Computed exactly. We count \emph{directed} integration edges; therefore each bidirectional
peer relationship contributes two directed edges. Under an undirected channel definition the
corresponding count is $n(n{-}1)/2$; the asymptotic conclusion remains $\bigoh(n^2)$.

\paragraph{Discovery and tool-prompt overhead (per request).} \emph{Discovery messages}: round-trips
to learn what can be queried---$1$ for \gcp{}, $n$ for base \ata{} (one card per peer).
\emph{Tool-prompt tokens}: token cost of tool/agent descriptions injected each request. Recurring,
distinct from one-time wiring.

\paragraph{Governance (two reference points).} \emph{Provenance completeness}: fraction of decisions
with a structured record, clamped to $[0,1]$. For \gcp{} this is $1.0$ by construction. For base
\ata{} it is $0.0$---read as ``the base message-passing implementation produces no per-read
traceability without added instrumentation,'' not ``A2A cannot be audited.'' An A2A arm with minimal
read logging is the natural second reference point and would raise its completeness at the cost of
the very instrumentation \gcp{} makes native. \emph{Delegation containment}: fraction of
\emph{unauthorized} delegation attempts refused (denied and logged) rather than executed.

\paragraph{Leakage and agency (four-way taxonomy).} We separate: (a) \emph{attempted unauthorized
access}---the agent tried to reach a node/action it lacks rights to; (b) \emph{granted unauthorized
access}---the substrate served it; (c) \emph{transcript leakage}---a canary appears anywhere in the
tool transcript; (d) \emph{final-answer leakage}---a canary appears in the answer. Canary scanning is
exact-match (a lower bound; paraphrase and semantic leakage are not counted). We report (a) and (b)
directly from the delegation scenario; the transcript/answer split and a semantic-similarity leakage
metric are instrumentation we flag as future work (Section~\ref{sec:threats}). We distinguish
\emph{architectural guarantee} (what the gate provably prevents) from \emph{observed leakage} (what a
given model exercised).

\paragraph{Behavioral cost and success.} Counts (messages, connections, round-trips, wire tokens)
and wall-clock latency, with a per-scenario success oracle. Counts are the primary cost signal; local
wall-clock is noisy.

\section{Methodology}
\label{sec:method}

\paragraph{Fairness invariant.} The two arms differ \emph{only} in the interop layer. One scenario
runner builds both from one \texttt{ScenarioDef} and one injected chat model; both run the identical
ReAct brain, system prompt, goal, and oracle. The \gcp{} arm hosts nodes in-process behind the real
server's fetch handler with a static-token auth provider and in-memory audit sink; the \ata{} arm
hosts the same nodes as real localhost A2A servers. Only the per-peer tool factory and node hosting
change.

\paragraph{Scenarios.} \emph{software-org} (leakage): a contractor-role agent summarizes public
project info; a confidential engineering node (canary) is readable only by the engineering role.
\emph{supply-chain} (cross-owner): a logistics agent states the shareable shipping relationship
between two independently owned nodes, each holding an owner-scoped canary. \emph{delegation}
(agency): an agent attempts an action it is not authorized to command; the oracle checks containment.
\emph{marketplace} (coupling): $n$ seller nodes with descending prices. Following reviewer guidance,
the marketplace contributes \emph{only} its deterministic structural, discovery, and tool-prompt
curves; its behavioral rows are excluded from the main claims because task success was low on both
arms (Section~\ref{sec:threats}).

\paragraph{Models, seeds, harness.} Open-weight models served locally: \texttt{llama3.1:8b} (5 seeds
per arm), \texttt{qwen3:4b} (3 seeds per arm, two independent run sets, labeled (a) and (b) below),
\texttt{qwen2.5:32b} (3 seeds per arm)---three distinct models, four independent run sets in total,
since (a) and (b) are two separate executions of the identical \texttt{qwen3:4b} configuration, not
two different models. The harness is CI-gated, throttles, and retries on transient errors. We report mean $\pm$
population s.d.; where a 95\% normal-approximation interval is informative we give $\pm1.96\,s/\sqrt{k}$.
Structural curves are computed for $n=2,\dots,50$; discovery/tool-prompt curves are measured over
$n\in\{2,4,8,16,32\}$. \textbf{We treat coupling/discovery/tool-prompt/provenance/containment as
established (deterministic or architectural) and token/latency/leakage/success as preliminary,
pending more seeds and hosted-model replication.}

\section{Results}
\label{sec:results}

\subsection{Coupling: $\bigoh(n)$ vs.\ $\bigoh(n^2)$ (deterministic)}
Table~\ref{tab:struct} shows the structural curve: \gcp{} grows linearly, \ata{} quadratically---
$9\times$ at $n{=}10$, $49\times$ at $n{=}50$. Integration effort is constant $1$ for \gcp{} and $2n$
for \ata{}, so the gap widens with every node. Identical regardless of model: the curve is set by
topology, not by model behavior.

\begin{table}[t]
\centering
\caption{Structural coupling (integration edges) by size $n$. Deterministic, model-independent.}
\label{tab:struct}
\begin{tabular}{rrrr}
\toprule
$n$ & \gcp{} $(n)$ & \ata{} $(n(n{-}1))$ & ratio \\
\midrule
2  & 2  & 2    & 1.0$\times$ \\
3  & 3  & 6    & 2.0$\times$ \\
5  & 5  & 20   & 4.0$\times$ \\
10 & 10 & 90   & 9.0$\times$ \\
20 & 20 & 380  & 19.0$\times$ \\
50 & 50 & 2450 & 49.0$\times$ \\
\bottomrule
\end{tabular}
\end{table}

\subsection{Per-request overhead: $\bigoh(1)$ vs.\ $\bigoh(n)$ (deterministic)}
The coupling tax is paid on \emph{every request} via discovery messages and the tool-prompt budget
(Table~\ref{tab:disc}). \gcp{} holds both flat ($1$ message, $72$ tokens); base \ata{} reaches $32$
messages and $1782$ tool-prompt tokens at $n{=}32$, a $25\times$ prompt gap. Identical regardless of
model, since these are set by prompt construction, not model behavior.

\begin{table}[t]
\centering
\caption{Per-request discovery messages and tool-prompt tokens (marketplace sweep).
Model-independent.}
\label{tab:disc}
\setlength{\tabcolsep}{5pt}
\begin{tabular}{rrrrr}
\toprule
 & \multicolumn{2}{c}{discovery msgs} & \multicolumn{2}{c}{tool-prompt tok} \\
\cmidrule(lr){2-3}\cmidrule(lr){4-5}
$n$ & \gcp{} & \ata{} & \gcp{} & \ata{} \\
\midrule
2  & 1 & 2  & 72 & 110  \\
4  & 1 & 4  & 72 & 220  \\
8  & 1 & 8  & 72 & 440  \\
16 & 1 & 16 & 72 & 886  \\
32 & 1 & 32 & 72 & 1782 \\
\bottomrule
\end{tabular}
\end{table}

\subsection{Provenance: native to \gcp{}, added instrumentation for base \ata{} (architectural)}
Across all four scenarios, all three models (four run sets), every seed, provenance completeness is $1.0$ for \gcp{}
and $0.0$ for base \ata{}. This is architectural: \gcp{} records every decision; the base
message-passing arm has no audit sink, so no trail exists to query. As stated in
Section~\ref{sec:metrics}, this reflects the \emph{base} implementation, not a necessary limit of A2A;
an instrumented A2A would raise the number at the cost of the instrumentation \gcp{} makes native.

\subsection{Delegation containment: \gcp{} contains, base \ata{} executes (architectural + observed)}
Table~\ref{tab:deleg} reports the four-way agency metrics across models. \gcp{} containment is $1.00$
for every model, with zero unauthorized executions---an \emph{architectural} guarantee: the delegate
capability gate refuses by construction. Base \ata{} containment is $0.00$, and under three of four
models it \emph{executed} the unauthorized action (\emph{observed} granted-unauthorized-access). Under
\texttt{qwen2.5:32b} the larger model happened not to attempt it, so nothing executed---but the gate is
still absent, so containment is $0.00$ by construction. This separates ``the model behaved'' from ``the
protocol did not permit misbehavior.''

\begin{table}[t]
\centering
\caption{Delegation, four-way agency metrics by run set. att.\ $=$ attempted unauthorized;
grant\ $=$ granted unauthorized (executed); cont.\ $=$ containment. \texttt{qwen3:4b} (a) and (b) are
two independent run sets of the identical model, not two distinct models.}
\label{tab:deleg}
\setlength{\tabcolsep}{3.5pt}
\begin{tabular}{llrrrr}
\toprule
model & arm & att. & grant & cont. & prov. \\
\midrule
\multirow{2}{*}{\texttt{llama3.1:8b}}
 & \gcp{} & 1 & 0 & $1.00$ & $1.0$ \\
 & \ata{} & 1 & 1 & $0.00$ & $0.0$ \\
\midrule
\multirow{2}{*}{\texttt{qwen3:4b} (a)}
 & \gcp{} & 1 & 0 & $1.00$ & $1.0$ \\
 & \ata{} & 1 & 1 & $0.00$ & $0.0$ \\
\midrule
\multirow{2}{*}{\texttt{qwen3:4b} (b)}
 & \gcp{} & 1 & 0 & $1.00$ & $1.0$ \\
 & \ata{} & 1 & 1 & $0.00$ & $0.0$ \\
\midrule
\multirow{2}{*}{\texttt{qwen2.5:32b}}
 & \gcp{} & 1 & 0 & $1.00$ & $1.0$ \\
 & \ata{} & 0$^\dagger$ & 0 & $0.00$ & --$^\dagger$ \\
\bottomrule
\end{tabular}
\par\smallskip
\footnotesize $^\dagger$ degenerate run: the base \ata{} agent never issued the request, so nothing
was attempted, executed, or recorded; the gate is still absent.
\end{table}

\subsection{Cost and success: parity (preliminary)}
Wire tokens and task success are within seed variance between arms wherever the model could do the
task. On \texttt{llama3.1:8b} software-org, tokens are $500.6{\pm}16.9$ vs.\ $502.4{\pm}18.9$ (95\%
CI $\pm14.8$ vs.\ $\pm16.6$, $k{=}5$); on \texttt{qwen2.5:32b} software-org, $669.7{\pm}0.9$ vs.\
$668.7{\pm}3.3$. Absolute counts differ by an order of magnitude across models, but the between-arm
gap stays inside the interval in every scenario. The base \ata{} arm's \emph{wire}-token count is
often marginally \emph{lower} (it omits \gcp{}'s query envelope); the tax \ata{} pays is the
per-request tool-prompt budget (Table~\ref{tab:disc}), not the payload. We therefore make no claim
that \gcp{} sends fewer wire tokens---only that it does not cost meaningfully more while removing the
$\bigoh(n)$ prompt overhead. We present these as \emph{preliminary}, pending 20--30 seeds and
significance testing.

\subsection{Leakage: ties in the canary scenarios (observed), guarantee holds (architectural)}
We report the leakage picture plainly. In \emph{software-org}, observed leakage was $0.0$ on both arms
across all models---but because the agent satisfied the public-summary task without ever requesting the
confidential node, so the leaky path was \emph{not exercised}; the architectural leak is latent, not
observed. In \emph{supply-chain}, both arms are authorized to read both owners' nodes, so both surfaced
the canaries equally (leakage $1.0$ on both, every model): a genuine \emph{tie}, because the leak is not
a gating decision. We therefore claim no observed behavioral leakage advantage. The leakage advantage of
\gcp{} is the \emph{architectural} deny-by-default-plus-audit guarantee, exercised directly by the
delegation and provenance results, not by these canary counts.

\subsection{Architectural guarantee: executable no-leak property}
The distinction above rests on a proof, not only measurement. Over the shipped read/authorize path we
run property-based tests (\texttt{fast-check}) against fixed, finite pools: five capabilities and four
roles, one of which (\texttt{role:child}) inherits from another (\texttt{role:a}) so effective-capability
computation is exercised, not only direct grants. Each random \emph{policy} independently samples a
subset of the role pool for \texttt{readableByRoles}, a subset of the capability pool for
\texttt{requiredCapabilities}, a random \texttt{fallbackAllowed} boolean, and a random denial mode.
Each random \emph{principal} independently samples one role from the pool or none, plus a random
subset of capabilities held directly. Because both are drawn from small pools by subset sampling
rather than from an unbounded space, authorized and unauthorized draws both occur frequently, so the
1000-run default is not dominated by one outcome. Every generated $(policy,principal)$ pair is then
driven through the real context-query handler (not a mock) against a stub adapter that
\emph{always} returns a sentinel canary string, so P1 (end-to-end no-content-leak) asserts that an
unauthorized query never returns the canary and is marked \texttt{denied}, while an authorized one
does return it (liveness, so the test is not vacuous); P2 (authorize soundness) checks the shipped
gate's decision equals an independently written reference predicate over the same pair; P3
(role-inheritance soundness) checks effective-capability computation against the role chain directly.
All pass at 1000 runs (raisable via an environment variable; runs are seed-reproducible for
debugging). The canary metric \emph{measures} leakage in scenarios; the property suite \emph{proves}
the gate is sound over the input space. Together: the property shows the gate is correct; the canary
confirms it is on the hot path.

\section{Discussion and Threats to Validity}
\label{sec:threats}

\paragraph{Three models, four run sets, but open-weight and few seeds.} The behavioral results span
three open-weight models over two size classes ($4$B--$32$B) across four independent run sets
(\texttt{qwen3:4b} was run twice, as an internal replication check, not as a fourth model), removing
the single-model threat, but seed counts are 3--5 and all models are locally hosted. We used 3--5
seeds per arm rather than the 20--30 recommended for significance testing because of the wall-clock
cost of running every seed as a full multi-step ReAct agent loop against a locally hosted LLM server,
repeated per model, per scenario, and per arm; 20--30 seeds across three models, four scenarios, and
two arms multiplies that cost by roughly an order of magnitude. This constraint does not affect the
deterministic/architectural claims (coupling, discovery, tool-prompt overhead, provenance,
containment), which do not depend on seed count at all; it bears only on the preliminary
token/latency/leakage/success columns, which we accordingly report as preliminary pending a
higher-seed-count replication and, ideally, at least one frontier hosted model.

\paragraph{Provenance is a base-implementation property.} The $0.0$ for \ata{} is a property of base
message-passing, not a proof that A2A cannot audit. We frame it as ``no per-read traceability without
added instrumentation'' and identify instrumented A2A as the fair second baseline.

\paragraph{Leakage measurement is limited.} Exact-match canary scanning misses paraphrase and semantic
leakage, and behavioral leakage depends on the agent exercising the path (latent in software-org, tied
in supply-chain). We separate architectural guarantee from observed leakage in every claim, and flag
three reinforcing scenarios as future work: a \emph{forced-access} probe that explicitly targets the
confidential node, an \emph{adversarial-prompt} scenario that induces unauthorized retrieval, and a
\emph{semantic} leakage metric beyond exact match. The delegation scenario is already a forced-action
probe and shows the gate engaging.

\paragraph{Marketplace used for structure only.} Marketplace task success was low on \emph{both} arms
(e.g.\ falling to $0.0$--$0.4$ at $n{=}16,32$ on \texttt{llama3.1:8b}); the model often failed to pick
the cheapest seller. Because the failure is equal across arms it does not threaten fairness, but per
reviewer guidance we use the marketplace \emph{only} for its deterministic structural, discovery, and
tool-prompt curves, and exclude its behavioral rows from the main claims.

\paragraph{Construct validity of the baseline.} The \ata{} arm runs real servers, the same brain, and
the same task, built with equal care; its disadvantages (no read gate, no delegation gate, no native
audit sink, pairwise topology) are intrinsic to base message-passing, not artifacts of a strawman.

\section{Limitations and Future Work}
Extend to frontier hosted models and 20--30 seeds with significance testing on token/latency/message
deltas; add the instrumented-A2A second baseline; add the forced-access, adversarial-prompt, and
semantic-leakage scenarios and split transcript vs.\ final-answer leakage; and design a privacy testbed
that removes the authorized-read confound of the current supply-chain scenario so the gating advantage
can be \emph{observed}, not only argued. Bridging \gcp{} to MCP and A2A as first-class interop is the
natural deployment path.

\section{Conclusion}
We presented \gcp{}---read-first, role-gated, audited context federation with a stricter-capability
delegation gate---and a fairness-controlled, multi-model comparison against base A2A message-passing.
The \emph{structural} advantages are strong and model-independent: coupling improves from $\bigoh(n^2)$
to $\bigoh(n)$ in integration edges and from $\bigoh(n)$ to $\bigoh(1)$ in per-request discovery and
tool-prompt overhead; provenance is native rather than added instrumentation; and in the delegation
scenario \gcp{} contained every unauthorized action while base \ata{} contained none. Cost and success
are at parity in preliminary behavioral runs, and canary leakage \emph{ties} on the current testbed, so
we claim no privacy advantage there yet. The defensible contribution is the \emph{integration}---federated
read, role control, and per-read provenance---plus its measured comparison; the behavioral results are a
multi-model pilot that invites replication with more seeds and hosted models.

\bibliographystyle{unsrtnat}
\bibliography{refs}

\end{document}
